\section*{अध्याय \ref{ch:b2:phylogenetic-trees} --- जातिवृत्तीय वृक्ष और आणविक घड़ी}

\begin{solution}{exo:b2:phylogenetic-trees:1}
\omterm{def:b2:phylogenetic-trees:tree}{समजातता}: किसी साझा पूर्वज से वंशागत कोई लक्षण --- किसी
चमगादड़ के पंख तथा किसी ह्वेल के फ़्लिपर की प्रगंडिका, बहिःप्रकोष्ठिका तथा अंतःप्रकोष्ठिका।
\omterm{def:b2:phylogenetic-trees:tree}{समवृत्तिता}: स्वतंत्र रूप से विकसित कोई एक जैसा लक्षण --- उड़ान-सतहों के रूप में
चमगादड़ों तथा पक्षियों के पंख (वे हड्डियाँ अग्र अंगों के रूप में समजात हैं,
पर पंख वैसे के वैसे नहीं), या ऑक्टोपस तथा
कशेरुकी की कैमरा-आँख। \omterm{def:b2:phylogenetic-trees:tree}{सहअपसारी लक्षण}: किसी \omterm{def:b2:phylogenetic-trees:tree}{अनुजात} में साझा और उसे परिभाषित करती
कोई अपसारी \omterm{def:b2:phylogenetic-trees:tree}{समजातता} --- पक्षियों के लिए पंख, उल्बधारियों के लिए
उल्बीय अंड।
\end{solution}

\begin{solution}{exo:b2:phylogenetic-trees:2}
पक्षी: एकवंशीय। पक्षियों के बिना सरीसृप: पारवंशीय। चतुष्पादों के बिना
मछली: पारवंशीय (फुफ्फुसमीन ट्राउट से अधिक हमारे पास हैं)।
उड़ने वाले कशेरुकी (चमगादड़, पक्षी, टेरोसॉर): बहुवंशीय।
स्तनी: एकवंशीय। प्रोकैरियोट: पारवंशीय (जीवाणु तथा
आर्किया, यूकैरियोटों को छोड़कर --- और आर्किया जीवाणुओं से अधिक
यूकैरियोटों के पास हैं)।
\end{solution}

\begin{solution}{exo:b2:phylogenetic-trees:3}
पक्षियों का सहोदर: मगर। स्तनियों का सहोदर: बाक़ी सारे
उल्बधारी (कछुए, छिपकलियाँ, मगर, पक्षी साथ)। छिपकलियों तथा पक्षियों को
रखता सबसे छोटा \omterm{def:b2:phylogenetic-trees:tree}{अनुजात}: (छिपकलियाँ,(मगर, पक्षी)),
यानी इस वृक्ष पर कछुओं के बिना द्विकोटरी। मगरों तथा
पक्षियों को उनकी गाँठ पर बदलने से कुछ नहीं बदलता: कोई गाँठ स्वतंत्र रूप से घुमाई जा सकती है।
\end{solution}

\begin{solution}{exo:b2:phylogenetic-trees:4}
अमूलित: $3$; $3\times 5\times 7 = 105$; $3\times 5\times
7\times 9\times 11\times 13\times 15 = \num{2027025}$। 4 वर्गकों के लिए
मूलित वृक्ष: तीनों अमूलित वृक्षों में से हर एक की 5 शाखाएँ हैं, यानी 15।
\end{solution}

\begin{solution}{exo:b2:phylogenetic-trees:5}
$((A,B),(C,D))$: स्थान 1, 2, 3 हर एक एक क़दम, स्थान 4 दो क़दम
(दोनों ओर $\mathrm{A}$ तथा $\mathrm{G}$), स्थान 5 एक: यानी 6 क़दम।
$((A,C),(B,D))$: $2 + 2 + 2 + 1 + 1 = 8$। $((A,D),(B,C))$: $2 + 2 + 2 +
2 + 1 = 9$। पहला वृक्ष सबसे मितव्ययी है; और उस पर स्थान 4
समरूप है ($\mathrm{A}\to\mathrm{G}$ परिवर्तन, या उसका उलटा,
दो बार होता है)।
\end{solution}

\begin{solution}{exo:b2:phylogenetic-trees:6}
$A$ तथा $B$ को ऊँचाई 3 पर जोड़िए। फिर $d(AB,C) = 14$, $d(AB,D) = 16$,
$d(C,D) = 10$: $C$ तथा $D$ को ऊँचाई 5 पर जोड़िए। फिर $d(AB,CD) = (14 + 16
+ 14 + 16)/4 = 15$: मूल ऊँचाई $7.5$ पर। वृक्ष $((A,B),(C,D))$।
\end{solution}

\begin{solution}{exo:b2:phylogenetic-trees:7}
$d = -0.75\ln(1 - 4p/3)$: $0.052$, $0.38$, $1.21$। $p = 0.6$ पर उस
लघुगणक का तर्क $0.2$ है और ढाल $\mathrm{d}d/\mathrm{d}p
= 1/(1 - 4p/3) = 5$: यानी $0.02$ की कोई प्रतिचयन-त्रुटि, जो $p$ में हो,
$0.1$ बन जाती है, यानी $d$ में; और हर स्थान पर बराबर दरों की मान्यता,
जो हर असली जीन में झूठी है, उसी आकार का कोई अभिनत जोड़ देती है। वह जीन
संतृप्त है।
\end{solution}

\begin{solution}{exo:b2:phylogenetic-trees:8}
$t = d/2k = 0.02/(2\times 10^{-9}) = 10$ मिलियन वर्ष। 500
मिलियन वर्ष पर, $d = 2\times 10^{-9}\times 5\times 10^{8} = 1.0$ और
$p = 0.75(1 - e^{-4/3}) = 0.55$: यानी संतृप्ति तक के रास्ते का तीन
चौथाई, जहाँ वह संशोधन तीखा है और वह तिथि भरोसे लायक नहीं
--- ऐसे विभाजनों के लिए कोई धीमा जीन चाहिए।
\end{solution}

\begin{solution}{exo:b2:phylogenetic-trees:9}
\qty{52}{\%}: वह \omterm{def:b2:phylogenetic-trees:tree}{अनुजात} मुश्किल से आधे पुनःप्रतिचयनों में फिर मिलता है;
वे आँकड़े उस पर लगभग चुप हैं, और वे विकल्प मिलकर
उतने ही संभाव्य हैं। \qty{99}{\%}: वह संकेत स्थानों भर
संगत है; वह \omterm{def:b2:phylogenetic-trees:tree}{अनुजात} ज़ोरदार ढंग से समर्थित है (उस प्रतिरूप तथा उस
पंक्तिबद्धन को देखते हुए --- \omterm{prop:b2:phylogenetic-trees:pitfalls}{लंबी-शाखा आकर्षण} जैसी व्यवस्थित त्रुटियाँ
बूटस्ट्रैप नहीं पकड़ता)। पहले के लिए, अनुक्रम जोड़िए (अधिक
जीन), उन शाखाओं को तोड़ने वाले वर्गक जोड़िए, और कोई भिन्न
विधि जाँचिए।
\end{solution}

\begin{solution}{exo:b2:phylogenetic-trees:10}
जो दो वंशक्रम हर एक स्थानों के किसी बड़े अंश पर बदल चुके हों
उनके पास, किसी भी ऐसे स्थान पर जहाँ दोनों बदले हों, उसी क्षार पर
उतरने की चार में एक संभावना है (बराबर दरों पर चार क्षारों के साथ)। यदि हर एक
स्थानों के \qty{40}{\%} पर बदला हो, तो वे दोनों
\qty{16}{\%} पर बदले हैं और सारे स्थानों के \qty{4}{\%} पर संयोग से मेल खाते हैं ---
जो उन स्थानों के अंश से बड़ा हो सकता है जो हर एक को उसके असली, धीरे विकसित होते
संबंधियों से जोड़ते सच्चे \omterm{def:b2:phylogenetic-trees:tree}{सहअपसारी लक्षण} ढोते हैं।
मितव्ययिता मेल गिनती है बिना यह पूछे कि वे क्यों होते हैं
और उन दोनों को जोड़ देती है। हर लंबी शाखा पर से निकलते वर्गक जोड़ना
उसे छोटे खंडों में बाँट देता है, ताकि वे
संयोग-मेल कई छोटी शाखाओं में बँट जाएँ और
हर खंड के सच्चे \omterm{def:b2:phylogenetic-trees:tree}{सहअपसारी लक्षण} दिखने लगें; और कोई
संभाव्यता-प्रतिरूप जो लंबी शाखाओं पर बहुत-से संयोग-मेलों की अपेक्षा करे
उनके लिए सीधे संशोधन कर देता है।
\end{solution}

\begin{solution}{exo:b2:phylogenetic-trees:11}
$d = -0.75\ln(1 - 0.12) = 0.096$; $t = d/2k = 0.096/(4\times 10^{-8})
= 2.4$ मिलियन वर्ष। प्रतिस्थापन: $0.096\times \num{16000} =
1500$, सापेक्ष परिशुद्धि $1/\sqrt{1500} = \qty{2.6}{\%}$: यानी
$\pm 0.06$ मिलियन वर्ष की कोई सांख्यिकीय त्रुटि। पर वह दर
दूसरे विभाजनों पर अंशांकित है और इस वंशक्रम के लिए दो के गुणक भर
ग़लत हो सकती है (वंशक्रमों के बीच दर-विविधता, और पीढ़ियों भर मापी गई
वंशावली-दर तथा करोड़ों वर्षों भर मापी गई दर के
बीच का अंतर); और उस जीन का
विचलन उन जातियों के विचलन से पहले का है, क्योंकि वह पैतृक
समष्टि बहुरूप थी --- इसलिए वह सच्चा विभाजन उस जीन के विभाजन से
हाल का है, और वह भी ऐसे काल भर जो उस पैतृक
समष्टि-आकार पर निर्भर है। (स्वीकृत तिथि, बहुत-से केंद्रकीय जीनों तथा
जीवाश्मों से, 6 से 7 मिलियन वर्ष है: यहाँ वह माइटोकॉन्ड्रियी घड़ी
बहुत तेज़ है।)
\end{solution}

\begin{solution}{exo:b2:phylogenetic-trees:12}
\omterm{prop:b2:phylogenetic-trees:pitfalls}{अपूर्ण वंशक्रम-छँटाई}: जब दो जाति-उद्भव काल में पास हों, तब
कोई पैतृक बहुरूपता ऐसे छँट सकती है कि किसी जीन का वृक्ष
जातियों को उस जाति-वृक्ष से भिन्न ढंग से समूहित करे; क्षैतिज
संचरण: किसी दूसरे वंशक्रम से पाया कोई जीन उस
वंशक्रम का इतिहास ढोता है; जीन-द्विगुणन: किसी एक जाति के किसी पैरालॉग की
किसी दूसरी के ऑर्थोलॉग से तुलना उस द्विगुणन की
तिथि देती है, उस जाति-उद्भव की नहीं। कोई जाति-वृक्ष
उस जीनोम भर से बहुत-से जीन लेकर और वह वृक्ष लेकर पाया जाता है जिसे बहुमत
समर्थन दे (या ऐसा कोई प्रतिरूप जो असंगत जीनों के किसी ज्ञात अंश की अपेक्षा करे),
ऑर्थोलॉग सावधानी से चुनकर, और
प्रोकैरियोटों में उन राइबोसोमी तथा सूचनात्मक जीनों को पसंद करके जो शायद ही कभी
संचरित होते हैं।
\end{solution}

\begin{solution}{pb:b2:phylogenetic-trees:1}
\textbf{1.} स्थान 1, 2 तथा 4 (ऐसी दशाएँ जो केवल $W$ तथा $H$ में साझा हैं,
और ऊँट के सापेक्ष अपसारी): यानी $\{W,H\}$ के \omterm{def:b2:phylogenetic-trees:tree}{सहअपसारी लक्षण}। स्थान 3
तथा 6: $\{W,H,C\}$ के \omterm{def:b2:phylogenetic-trees:tree}{सहअपसारी लक्षण}। स्थान 5 $W$ का कोई स्वअपसारी लक्षण है,
यानी किसी एक वंशक्रम में कोई परिवर्तन जो कुछ भी समूहित नहीं करता; और उस आणविक वृक्ष पर कोई स्थान
समरूप नहीं है।
\textbf{2.} प्रति स्थान एक परिवर्तन: यानी 6 क़दम।
\textbf{3.} स्थान 1, 2 तथा 4 को अब हर एक को दो परिवर्तन चाहिए, और स्थान 3, 5
तथा 6 को एक: यानी 9 क़दम।
\textbf{4.} ह्वेल--दरियाई घोड़ा वृक्ष, तीन क़दमों से। हज़ार में से
तीन स्थान कोई छोटा अंतर है जिसे कुछ समरूपताएँ पलट सकती हैं;
उसे मज़बूत करने का अर्थ है अधिक जीन, अधिक वर्गक (दूसरे
रोमंथी, पेकरी, हिरन), और कोई बूटस्ट्रैप।
\textbf{5.} वह \omterm{def:b2:phylogenetic-trees:tree}{अनुजात} छह में से पाँच पुनःप्रतिचयनों में लौटता है: यानी मध्यम
समर्थन, संयोग से कहीं ऊपर पर उस \qty{95}{\%} से कम
जिसे परंपरा से मज़बूत कहा जाता है।
\textbf{6.} वह बाह्यसमूह उस समूह के बाहर होना चाहिए पर इतना पास कि
उसका अनुक्रम अब भी पंक्तिबद्ध हो सके तथा असंतृप्त हो; और ऊँट
गाय--सूअर--दरियाई घोड़ा--ह्वेल समूह के बाहर कोई सम-खुरीय है। कोई
मछली लगभग संतृप्त स्तरों पर भिन्न होती, उसकी दशाएँ स्तनियों के भीतर इस बारे में
सूचना न देतीं कि कौन-सी दशा पैतृक है,
और उसकी लंबी शाखा सबसे तेज़ विकसित होते अंतःसमूह वंशक्रम को
आकर्षित कर लेती।
\textbf{7.} $d = 0.0625$, $0.107$, $0.131$, $0.155$।
\textbf{8.} $W$ तथा $H$ को ऊँचाई $0.031$ पर जोड़िए। फिर $d(WH,C) =
0.107$, $d(WH,P) = 0.131$, $d(C,P) = 0.131$, और $M$ तक सब कुछ
$0.155$: $WH$ तथा $C$ को $0.054$ पर जोड़िए। फिर $d(WHC,P) = 0.131$: $0.065$
पर जोड़िए। और अंत में $M$ को $0.078$ पर। वृक्ष $(M,(P,(C,(W,H))))$।
\textbf{9.} हाँ: वही वृक्ष, वही नेस्टिंग।
\textbf{10.} जो ह्वेल दुगुना तेज़ बदली होती वह
दरियाई घोड़े समेत हर चीज़ से दूर होती; और UPGMA, जो हर
सिरा ऊँचाई शून्य पर रखती है और सबसे पास के जोड़े को जोड़ती है, पहले दरियाई घोड़े को
गाय से जोड़ देती और उस ह्वेल को बाहर छोड़ देती --- यानी ठीक
शरीर-रचनाविदों का वृक्ष, पर ग़लत कारण से।
\textbf{11.} $0.0625\times 1000 = 62.5$ प्रतिस्थापन; पॉइसों
मानक विचलन $\sqrt{62.5} = 7.9$।
\textbf{12.} $7.9/62.5 = \qty{13}{\%}$।
\textbf{13.} $k = d/2t = 0.107/(1.2\times 10^{8}) = 8.9\times 10^{-10}$
प्रति स्थान प्रति वर्ष।
\textbf{14.} $t = 0.0625/(2\times 8.9\times 10^{-10}) = 35$ मिलियन
वर्ष।
\textbf{15.} सूअर: $0.131/(1.79\times 10^{-9}) = 73$ मिलियन वर्ष;
ऊँट: $0.155/(1.79\times 10^{-9}) = 87$ मिलियन वर्ष।
\textbf{16.} $35 \pm 4.5$ मिलियन वर्ष (\qty{13}{\%})।
\textbf{17.} वह दर $1/t_{\text{अंश}}$ की तरह चलती है और वह तिथि
$t_{\text{अंश}}$ की तरह: $\pm 5/60 = \qty{8}{\%}$, इसलिए $\pm 3$ मिलियन वर्ष
--- और उस पॉइसों त्रुटि से मिलाकर, लगभग $\pm 5$ मिलियन वर्ष।
\textbf{18.} $d = -0.75\ln(1 - 0.6) = 0.69$: प्रभाव में वह जीन अपने
स्थानों के दो तिहाई पर बदल चुका है, उस संशोधन ने
$p$ को $1.5$ से गुणा किया है और उसकी ढाल $2.5$ है; वह इस
विभाजन के लिए संतृप्ति के पास है और उसकी तिथि थोड़े ही लायक़ है।
\textbf{19.} किसी लक्षण की अनुपस्थिति कोई साझा अपसारी दशा नहीं है:
ह्वेलों ने पश्च अंग तथा उसका टख़ना पूरा खो दिया, इसलिए इस बारे में कुछ नहीं
कहा जा सकता कि उनका कौन-सा टख़ना होता। कोई \omterm{def:b2:phylogenetic-trees:tree}{सहअपसारी लक्षण}
किसी \omterm{def:b2:phylogenetic-trees:tree}{अनुजात} का प्रमाण है; किसी एक वंशक्रम में उसका खोना उस सदस्यता के
विरुद्ध प्रमाण नहीं है।
\textbf{20.} उन जीवाश्म ह्वेलों में दुहरी-घिरनी टख़ना है: यानी ह्वेल
ऐसे किसी पूर्वज से उतरी हैं जिसके पास वह था, और वही उन्हें
सम-खुरीयों के भीतर रखता है। वह आणविक भविष्यवाणी उन
चट्टानों से पुष्ट होती है।
\textbf{21.} पारवंशीय: यानी कोई पूर्वज अपने सारे वंशजों समेत,
सिवाय ह्वेल-रेखा के।
\textbf{22.} दो तेज़ जाति-उद्भवों के आरपार किसी पैतृक बहुरूपता की
\omterm{prop:b2:phylogenetic-trees:pitfalls}{अपूर्ण वंशक्रम-छँटाई} (गाय, दरियाई घोड़ा तथा ह्वेल कुछ ही
मिलियन वर्षों के भीतर अलग हुए), या पैरालॉगों की तुलना। निर्णय
बहुत-से जीनों का अनुक्रमण करके और वह वृक्ष लेकर कीजिए जिसका उनमें से अधिकांश, तथा
वह संयोजन, समर्थन करें; और हर कुल में द्विगुणन की जाँच कीजिए।
\textbf{23.} $d = -0.75\ln(1 - 0.333) = 0.30$; $k = d/2t =
0.30/(7\times 10^{7}) = 4.3\times 10^{-9}$ प्रति स्थान प्रति वर्ष, यानी उस
केंद्रकीय जीन से पाँच गुना।
\textbf{24.} माइटोकॉन्ड्रियी DNA का प्रतिकरण ऐसी पॉलिमरेज़ करती है जिसकी
शुद्धिपाठन कमज़ोर है, वह श्वसन-शृंखला के ऑक्सीजन-मूलकों के सामने पड़ा है,
और उसकी मरम्मत केंद्रकीय मशीनरी नहीं करती; और उसमें
पुनर्योजन भी नहीं है, इसलिए क्षति जमा होती जाती है।
\textbf{25.} वृक्ष $(M,(P,(C,(W,H))))$: यानी ह्वेल दरियाई घोड़ों का सहोदर
समूह हैं; $k = 8.9\times 10^{-10}$ प्रति स्थान प्रति वर्ष; और ह्वेल--दरियाई घोड़ा
विचलन $35 \pm 5$ मिलियन वर्ष।
\end{solution}
